\paragraph{\trFunctoriality。}\label{sub:functoriality}
设 $\scr{G}\colon\scr{A}\to\scr{B}$ 是一个\trNonunital{} \trAFunctor。诱导\trFunctor{}%
$\Sigma\scr{G}\colon\Sigma\scr{A}\to\Sigma\scr{B}$ 定义为
\begin{equation}\label{eq:sigma-functor}
\begin{aligned}
&(\Sigma\scr{G})\Bigl(\bigoplus_{i\in I}V^i\otimes X^i\Bigr)
=\bigoplus_{i\in I}V^i\otimes\scr{G}(X^i),\\
&(\Sigma\scr{G})^d(a_d,\ldots,a_1)^{i_d,i_0}
=\sum_{i_1,\ldots,i_{d-1}}(-1)^{\lhd}\,
\phi_d^{i_d,i_{d-1}}\cdots\phi_1^{i_1,i_0}\otimes\\
&\hspace{6em}\otimes\,\scr{G}^d\bigl(x_d^{i_d,i_{d-1}},\ldots,x_1^{i_1,i_0}\bigr),
\end{aligned}
\end{equation}
其中 $\lhd$ 与诸 $a_k$ 均同\eqref{eq:sigma-composition}。此外，任一
$T\in\hom_{\mathit{nu}\text{-}\mathit{fun}(\scr{A},\scr{B})}(\scr{G}_0,\scr{G}_1)$
都诱导一个
$\Sigma^1T\in\hom_{\mathit{nu}\text{-}\mathit{fun}(\Sigma\scr{A},\Sigma\scr{B})}
(\Sigma\scr{G}_0,\Sigma\scr{G}_1)$，由下式类似地给出：
\begin{equation}\label{eq:sigma1-map}
\begin{split}
(\Sigma^1T)^d(a_d,\ldots,a_1)^{i_d,i_0}
&=\sum_{i_1,\ldots,i_{d-1}}(-1)^{\lhd}\,
\phi_d^{i_d,i_{d-1}}\cdots\phi_1^{i_1,i_0}\otimes\\
&\quad\otimes T^d\bigl(x_d^{i_d,i_{d-1}},\ldots,x_1^{i_1,i_0}\bigr).
\end{split}
\end{equation}
公式\eqref{eq:sigma-functor} 与\eqref{eq:sigma1-map}（高阶项取零）定义出一个%
\trNonunital{} \trAFunctor{}%
$\Sigma\colon\mathit{nu}\text{-}\mathit{fun}(\scr{A},\scr{B})\to
\mathit{nu}\text{-}\mathit{fun}(\Sigma\scr{A},\Sigma\scr{B})$。

\trTwistedComplex{}的情形完全平行。存在一个诱导的\trNonunital{} \trAFunctor{}%
$Tw\scr{G}\colon Tw\scr{A}\to Tw\scr{B}$：
\begin{equation*}
(Tw\scr{G})(X,\delta_X)
=\Bigl(\Sigma\scr{G}(X),\ \sum_e(\Sigma\scr{G})^e(\delta_X,\ldots,\delta_X)\Bigr),
\end{equation*}
\begin{equation*}
\begin{split}
(Tw\scr{G})^d(a_d,\ldots,a_1)
&=\sum_{i_0,\ldots,i_d}(\Sigma\scr{G})^{d+i_0+\cdots+i_d}\bigl(
\underbrace{\delta_{X_d},\ldots,\delta_{X_d}}^{i_d},a_d,\\
&\qquad \underbrace{\delta_{X_{d-1}},\ldots,\delta_{X_{d-1}}}_{i_{d-1}},
a_{d-1},\ldots,a_1,
\underbrace{\delta_{X_0},\ldots,\delta_{X_0}}_{i_0}\bigr).
\end{split}
\end{equation*}
注意，若 $X$ 带有由子\trComplex{} $F^kX$ 组成的\trFiltration，使得其\trDifferential{}成为
严格下三角的，那么由诸 $\Sigma\scr{G}(F^kX)$ 给出的 $\Sigma\scr{G}(X)$ 的诱导%
\trFiltration{}亦是如此，因此它确是一个\trTwistedComplex。对
$T\in\hom_{\mathit{nu}\text{-}\mathit{fun}(\scr{A},\scr{B})}(\scr{G}_0,\scr{G}_1)$，
定义
$Tw^1T\in\hom_{\mathit{nu}\text{-}\mathit{fun}(Tw\scr{A},Tw\scr{B})}
(Tw\scr{G}_0,Tw\scr{G}_1)$ 为
\begin{equation}\label{eq:tw1-map}
\begin{split}
(Tw^1T)^d(a_d,\ldots,a_1)
&=\sum_{i_0,\ldots,i_d}(\Sigma T)^{d+i_0+\cdots+i_d}\bigl(
\underbrace{\delta_{X_d},\ldots,\delta_{X_d}}^{i_d},a_d,\\
&\qquad \underbrace{\delta_{X_{d-1}},\ldots,\delta_{X_{d-1}}}_{i_{d-1}},
a_{d-1},\ldots,a_1,
\underbrace{\delta_{X_0},\ldots,\delta_{X_0}}_{i_0}\bigr).
\end{split}
\end{equation}
\begingroup\emergencystretch=8em\hbadness=10000
如前，这一切合起来构成一个\trNonunital{} \trAFunctor{}%
$Tw\colon\mathit{nu}\text{-}\mathit{fun}(\scr{A},\scr{B})\to$
$\mathit{nu}\text{-}\mathit{fun}(Tw\scr{A},Tw\scr{B})$，并且它与左\trComposition{}和
右\trComposition{}\trCompatible。这意味着，对一个固定的 $\scr{G}\colon\scr{A}\to\scr{B}$，有
\par\endgroup
\begin{equation}\label{eq:tw-left-right-composition}
\begin{aligned}
\scr{L}_{Tw\scr{G}}\circ Tw=Tw\circ\scr{L}_{\scr{G}}\colon
\mathit{nu}\text{-}\mathit{fun}(\scr{C},\scr{A})
&\longrightarrow\mathit{nu}\text{-}\mathit{fun}(Tw\scr{C},Tw\scr{B}),\\
\scr{R}_{Tw\scr{G}}\circ Tw=Tw\circ\scr{R}_{\scr{G}}\colon
\mathit{nu}\text{-}\mathit{fun}(\scr{B},\scr{C})
&\longrightarrow\mathit{nu}\text{-}\mathit{fun}(Tw\scr{A},Tw\scr{C}).
\end{aligned}
\end{equation}

\vkthmsync{23}
\begin{lemma}\label{thm:tw-preserves-cohomologically-ff}
若 $\scr{G}\colon\scr{A}\to\scr{B}$ 是\trCohomology\trFullyFaithful{}的，那么
$Tw\scr{G}\colon Tw\scr{A}\to Tw\scr{B}$ 亦然。
\end{lemma}
\begin{proof}
对 $(Tw\scr{G})^1$ 在\trSpectralSequence{} \eqref{eq:twisted-spectral-sequence}
上的作用使用一个比较论证，问题便化归为关于 $\Sigma\scr{G}$ 的相应陈述，而后者是平凡的。
\end{proof}
